\section{Introduction}\label{sec:introduction}
A scheduler ranking is a conditional statement. It compares a fixed set of
implementations under a workload, an operating region, an objective and an
execution model. A ranking obtained from one trace cannot automatically be
used to select a policy for another trace; nor does agreement on average
throughput imply agreement on deadline attainment or tail latency. This
paper studies how such comparative conclusions transfer, rather than
introducing a new scheduler.

LLM serving makes this question concrete. Continuous batching and paged
key--value (KV) memory management permit concurrent requests
\citep{kwon2023vllm,yu2022orca}. Chunked prefill changes competition between
prompt processing and decoding \citep{agrawal2024sarathi}; disaggregation
separates their resource requirements \citep{zhong2024distserve,patel2024splitwise}.
Production traces differ in arrival bursts, prompt lengths and output lengths
\citep{burstgpt2025,stojkovic2025dynamollm}. These established observations
motivate testing the stability of a comparison across conditions. They do
not establish which policies will reverse order, or even that reversals
must occur in a particular evaluation.

We study the LLM-Serving Scheduler Portability Benchmark (LSSP), a frozen,
trace-driven simulator campaign. Its primary corpus contains 120 windows
from three separately released trace sources. A 13-policy execution panel,
including 11 primary policies and two style approximations, is evaluated
in six FIFO-referenced operating regions. There are 9,360 distinct
configurations and two deterministic verification passes per configuration.
The passes check repeatability of the simulator; they are not independent
samples. Uncertainty is instead assessed across workload windows, subject
to the dependence limitations in Section~\ref{sec:threats}.

The contribution is a reproducible evaluation methodology and an empirical
characterization of \emph{simulated} scheduler-ranking portability. The
methodology combines a fixed comparison panel, matched policy outcomes,
explicit metric eligibility, rank agreement, practical-margin reversal
classification and a finite-corpus ranking-recovery analysis. Kendall
correlation, bootstrap intervals and false-discovery correction are
standard statistical tools, not a new algorithm or theoretical result.
The scientific object is the reliability and scope of the comparative
conclusion produced by their application.

The primary ANWG ranking correlations range from 0.551 to 1.000 across 18
source-pair/region comparisons. Of 990 pairwise-policy comparison records,
36 meet the saved practical-margin and confidence-interval reversal criteria;
these span five policy pairs and all involve BurstGPT. Most comparisons
therefore retain their order. A separately labeled post-campaign analysis
shows a median cross-metric Kendall correlation of 0.419 and agreement of
the deterministic top-choice rule in 316 of 990 conditions. Agreement
and disagreement are both findings; neither is assumed in advance.

These are observations of historical simulator implementations. Later
source and engineering audits identified limits in partial-prefill
execution, native admission, KV allocation, preemption and execution
acknowledgment. The historical GPU comparison also has metadata and replay
limitations. It cannot establish faithful native-policy ranking transfer.
Accordingly, this paper neither claims physical validation of its simulated
reversals nor generalizes them to GPU architectures or parallel deployments.
The distinction is central to the domain of validity, not a qualification
added after interpreting the numerical results.

The original six research questions are retained in
Section~\ref{sec:rq}. Sections~\ref{sec:execution-model} and
\ref{sec:workloads} define the model and experimental design;
Sections~\ref{sec:within}, \ref{sec:results} and
\ref{sec:results-cross-metric} separate operating-region behavior,
cross-source portability and cross-metric portability. Robustness and
practical implications precede an explicit account of threats to validity.
